\chapter{Motifs intrajournaliers}\label{ch:s1:intraday-patterns}

Le marché actions américain se négocie de 9~h~30 à 16~h, mais ses contrats à terme se négocient vingt-quatre heures sur vingt-quatre, et Bondarenko et Muravyev ont établi que quatre heures de la nuit,
autour de l'ouverture européenne, expliquent à elles seules l'ensemble du rendement moyen de l'E-mini S\&P 500~; les vingt autres heures ont en moyenne un rendement nul, au bruit près. La journée de négociation a
sa propre forme~: le volume s'accumule à l'ouverture et à la clôture, la fixation du cours de clôture établit les prix auxquels les \omterm{def:s1:index-rebalancing:fund}{fonds indiciels} et les dealers valorisent leurs portefeuilles, et, dans la dernière
demi-heure, le marché tend à prolonger le mouvement de la journée, parce que les dealers vendeurs de gamma doivent couvrir dans son sens. Ce chapitre décompose la
journée, simule le flux de couverture et le déséquilibre à la clôture, et mesure ce que vaut chaque régularité pour la négociation. La construction est \texttt{firm.intraday}.

\section{Rendements overnight et rendements intrajournaliers}

\begin{definition}[Rendement overnight, rendement intrajournalier]\label{def:s1:intraday-patterns:returns}
Le \emph{rendement overnight}\index{rendement overnight} d'un titre est son rendement entre la clôture d'une journée et l'ouverture de la suivante~; son \emph{rendement
intrajournalier}\index{rendement intrajournalier} est le rendement entre l'ouverture et la clôture. Leur somme (en logarithmes) est le rendement clôture à clôture.
\end{definition}

La décomposition est étudiée depuis les premières données intrajournalières~: l'article «~like night and day~» de Cooper, Cliff et Gulen en est la référence classique, et Lou, Polk et
Skouras ont montré que les profits du retournement et de plusieurs stratégies de momentum étaient réalisés entièrement pendant la nuit, avec une continuation à l'intérieur de chaque partie de la journée
et un retournement d'une partie à l'autre (chapitre~2). La part du rendement propre au marché qui est gagnée la nuit est plus difficile à déterminer qu'il n'y paraît. \texttt{firm.intraday}
simule vingt ans d'un indice en barres d'une demi-heure, avec un rendement espéré de 7~\% réparti proportionnellement à la variance, un quart de la variance arrivant
pendant la nuit, et donc un quart du rendement espéré~: 1.75~\% par an la nuit. Mesurée sur les vingt ans, la part nocturne du rendement cumulé s'élève à
54~\% (4.4~\% par an la nuit, 3.7~\% le jour). L'erreur type d'un rendement moyen sur vingt ans, avec une volatilité de 16~\%, vaut 3.6~\% par an, et celle de sa composante
nocturne, 1.8~\%~: vingt ans ne permettent pas de distinguer un quart d'une moitié. Le résultat de Bondarenko et Muravyev est plus net, car il repose sur l'horloge complète des contrats à terme et
sur un ratio de Sharpe de 1.6 en quatre heures, et non sur la répartition d'une moyenne bruitée.

\section{Le smile du volume et les enchères}

\begin{definition}[Smile du volume]\label{def:s1:intraday-patterns:smile}
Le \emph{smile du volume}\index{smile du volume} est la forme en U du volume des échanges (et de la volatilité) au cours de la journée de négociation~: élevé à l'ouverture, faible en milieu de
journée, maximal à la clôture, la fixation du cours de clôture concentrant une part importante des échanges.
\end{definition}

Dans le modèle de ce chapitre, la première et la dernière demi-heure portent chacune 12.1~\% du volume de la journée, et la demi-heure de milieu de journée 4.9~\%~; un quart du volume
s'échange dans la première et la dernière demi-heure réunies. La clôture est le moment où sont fixés les prix de référence~: les \omterm{def:s1:index-rebalancing:fund}{fonds indiciels} (chapitre~10), les fonds communs qui valorisent leurs
parts à la clôture, et les dealers qui valorisent leurs portefeuilles négocient tous à ce moment-là, par la fixation du cours de clôture (livre~1, chapitre~13).

\begin{dated}{2026-09}{Le cross de clôture de Nasdaq}\label{dat:s1:intraday-patterns:nasdaq}
Dans la description que Nasdaq donne de son cross de clôture (édition 2019 de sa FAQ), les ordres au marché à la clôture, les ordres limite à la clôture et les ordres de déséquilibre seul sont acceptés avant 15~h~50
heure de l'Est~; à partir de 15~h~50, Nasdaq diffuse son indicateur de déséquilibre net des ordres (toutes les 10~secondes, puis toutes les secondes à partir de 15~h~55), et les ordres à la clôture
ne peuvent plus être annulés ni modifiés~; la saisie des ordres au marché à la clôture s'arrête à 15~h~55, celle des ordres limite à la clôture à 15~h~58, et les ordres de déséquilibre seul sont
acceptés jusqu'à 16~h~00, heure à laquelle commence le processus de clôture. Les autres bourses organisent leurs propres enchères, avec leurs propres heures limites.
\end{dated}

\begin{omfigure}
\begin{tikzpicture}
\begin{axis}[name=a,ybar,bar width=7pt,width=6.4cm,height=5.2cm,xlabel={\small demi-heure de la journée},ylabel={\small part du volume (\%)},
  tick label style={font=\footnotesize},xtick={1,4,7,10,13},ymin=0,ymax=14,grid=major,grid style={gray!25},table/col sep=comma]
\addplot[fill=omDef!60,draw=omDef] table[x=bar,y=share_pct]{figdata/strategies-1/13-intraday-patterns/smile.csv};
\end{axis}
\begin{axis}[at={(a.east)},anchor=west,xshift=1.3cm,width=6.8cm,height=5.2cm,xlabel={\small rendement avant la dernière demi-heure (\%)},
  ylabel={\small dernière demi-heure (pb)},tick label style={font=\footnotesize},grid=major,grid style={gray!25},
  legend style={font=\scriptsize,at={(0.5,-0.3)},anchor=north},table/col sep=comma]
\addplot[black!50,thick,mark=o] table[x=rest0,y=last0_bp]{figdata/strategies-1/13-intraday-patterns/binned.csv};
\addplot[omThm,very thick,mark=*] table[x=rest5,y=last5_bp]{figdata/strategies-1/13-intraday-patterns/binned.csv};
\legend{absence de flux de couverture,couverture des vendeurs de gamma}
\end{axis}
\end{tikzpicture}

\omcaption{Le modèle de journée du chapitre. À gauche~: part du volume quotidien par demi-heure. À droite~: rendement moyen de la dernière demi-heure par décile du rendement entre la clôture précédente et la
dernière demi-heure, sans et avec une couverture par des dealers vendeurs de gamma. Données~: \texttt{s1\_intraday}.}\label{fig:s1:intraday-patterns:smile}
\end{omfigure}

\section{Flux de couverture de fin de journée}

\begin{definition}[Flux de couverture de fin de journée]\label{def:s1:intraday-patterns:hedging}
Un \emph{flux de couverture de fin de journée}\index{flux de couverture de fin de journée} est la négociation que les dealers et les fonds à effet de levier ou inverses doivent effectuer près de la clôture pour rééquilibrer
leur exposition après le mouvement de la journée~; lorsqu'ils sont vendeurs de gamma (livre~1, chapitre~26), la couverture achète après les hausses et vend après les baisses, dans le sens du
mouvement.
\end{definition}

Baltussen, Da, Lammers et Martens ont mis en évidence un momentum intrajournalier sur plus de soixante contrats à terme sur actions, obligations, matières premières et devises, de 1974 à 2020~: le
rendement des trente dernières minutes avant la clôture est prédit positivement par le rendement entre la clôture précédente et ce moment, de façon significative et dans toutes
les classes d'actifs, et l'effet s'inverse au cours des jours suivants. Ils l'ont relié à la couverture en gamma des teneurs de marché sur options et des ETF à effet de levier. Le
modèle du chapitre introduit ce mécanisme~: les dealers vendeurs de gamma ajoutent à la dernière demi-heure 5~\% du mouvement de la journée jusqu'alors, et la moitié de la poussée s'inverse
pendant la nuit.

\begin{center}
\small
\begin{tabular}{@{}p{6.4cm}>{\centering\arraybackslash}p{3.6cm}>{\centering\arraybackslash}p{3.6cm}@{}}
\toprule
vingt ans synthétiques & sans flux de couverture & couverture vendeuse de gamma\\
\midrule
pente de la dernière demi-heure sur le reste de la journée~; $t$ & $-0.006$~; $-1.2$ & 0.044~; 9.2\\
momentum de la dernière demi-heure~: taux de réussite (hit rate)~; ratio de Sharpe après coûts & 48.8~\%~; $-0.77$ & 54.1~\%~; 1.05\\
momentum de la dernière demi-heure~: rendement par an & $-4.0~\%$ & 5.5~\%\\
contre-tendance sur le déséquilibre à la clôture~: ratio de Sharpe~; rendement par an & 1.09~; 8.8~\% & 1.08~; 8.7~\%\\
\omterm{def:s1:intraday-patterns:returns}{rendement overnight} suivant sur la dernière demi-heure (pente) & $-0.048$ & $-0.053$\\
\bottomrule
\end{tabular}
\end{center}

Sans flux de couverture, la dernière demi-heure est imprévisible (pente $-0.006$, $t = -1.2$), et une transaction qui suit la direction de la journée pendant la dernière demi-heure
perd, à cause de ses coûts d'un point de base par jour, $-4.0~\%$ par an. Avec ce flux, la pente vaut 0.044 ($t = 9.2$), la direction est juste dans 54.1~\% des cas, et la transaction
rapporte 5.5~\% par an, avec un ratio de Sharpe de 1.05 (la \cref{fig:s1:intraday-patterns:smile}, à droite). La transaction ne vaut que par la demande de couverture qui la sous-tend~:
lorsque les dealers sur options sont acheteurs de gamma (après que des investisseurs leur ont vendu des options), la couverture s'oppose au mouvement et la régularité s'inverse. Mesurer le
gamma net des dealers est la version praticienne du paramètre \texttt{gamma} de ce chapitre.

\section{Déséquilibres à la clôture}

La fixation du cours de clôture publie son déséquilibre avant de fixer le prix. Dans le modèle, un déséquilibre d'un écart type déplace la clôture de dix points de base, et
60~\% de cette poussée s'inverse pendant la nuit. Une transaction qui prend le côté opposé du déséquilibre publié à la clôture et sort à l'ouverture suivante rapporte 8.8~\% par an
avec un ratio de Sharpe de 1.09. La logique est celle du chapitre~10 en version réduite~: celui qui doit négocier à la clôture paie celui qui peut attendre l'ouverture. Le risque est lui aussi
le même~: un déséquilibre peut porter une information, et la part qui ne s'inverse pas constitue la perte sur les transactions qui n'auraient pas dû être réalisées.

\section{Fiches de stratégie}

\begin{strategyfile}{Dérive overnight}\label{strat:s1:intraday-patterns:overnight}
\sfield{Qui paie, et pourquoi} Des investisseurs qui détiennent du risque pendant la nuit, lorsque l'incertitude se dissipe autour de l'ouverture européenne, et qui sont rémunérés pour cela.
\sfield{Instruments et plateformes} Des contrats à terme sur indices négociés vingt-quatre heures sur vingt-quatre.
\sfield{Signal} L'heure de la journée~: position longue dans les heures autour de l'ouverture européenne.
\sfield{Dimensionnement et exécution} Entrer avant la fenêtre, sortir après elle~; positions petites et répétées.
\sfield{Coûts} Deux transactions sur contrats à terme par jour.
\sfield{Comment elle disparaît} En tant que prime de risque, elle se paie par des pertes les mauvaises nuits~; et par l'encombrement (crowding) sur la même fenêtre.
\sfield{Horizon, capacité, infrastructure} Quelques heures~; un desk de contrats à terme ouvert vingt-quatre heures sur vingt-quatre.
\sfield{Backtest honnête} Prix des contrats à terme aux heures exactes~; coûts~; la période COVID dans et hors de l'échantillon.
\sfield{Sources} Bondarenko et Muravyev (2022)~: l'ensemble du rendement moyen en quatre heures autour de l'ouverture européenne, ratio de Sharpe de 1.6.
\end{strategyfile}

\begin{strategyfile}{Momentum de la dernière demi-heure}\label{strat:s1:intraday-patterns:lasthalf}
\sfield{Qui paie, et pourquoi} Des dealers et des fonds à effet de levier qui doivent se couvrir dans le sens du mouvement de la journée.
\sfield{Instruments et plateformes} Des contrats à terme sur indices et d'autres contrats liquides.
\sfield{Signal} Le rendement depuis la clôture précédente jusqu'à 30~minutes avant la clôture.
\sfield{Dimensionnement et exécution} Dans le sens de ce rendement pour la dernière demi-heure~; position nulle à la clôture.
\sfield{Coûts} Un aller-retour par jour sur les contrats les plus liquides.
\sfield{Comment elle disparaît} Lorsque les dealers sont longs en gamma~; lorsque d'autres opérateurs négocient la même demi-heure en premier.
\sfield{Horizon, capacité, infrastructure} Une demi-heure~; données intrajournalières et exécution.
\sfield{Backtest honnête} Prix aux bornes exactes des barres~; coûts~; le retournement du lendemain.
\sfield{Sources} Baltussen, Da, Lammers et Martens (2021)~; Gao, Han, Li et Zhou (2018, chapitre~5).
\end{strategyfile}

\begin{strategyfile}{Retournement du déséquilibre à la clôture}\label{strat:s1:intraday-patterns:imbalance}
\sfield{Qui paie, et pourquoi} Les traders qui doivent négocier à la clôture et en poussent le cours.
\sfield{Instruments et plateformes} Les actions et leurs enchères de clôture.
\sfield{Signal} Le déséquilibre publié avant l'enchère.
\sfield{Dimensionnement et exécution} Compenser le déséquilibre dans l'enchère~; sortir à l'ouverture suivante ou au cours de la journée suivante.
\sfield{Coûts} Frais d'enchère~; risque overnight.
\sfield{Comment elle disparaît} Déséquilibres qui portent de l'information~; concurrence pour les compenser.
\sfield{Horizon, capacité, infrastructure} Overnight~; flux de données de déséquilibre d'enchère et heures limites de saisie des ordres.
\sfield{Backtest honnête} Le déséquilibre tel que publié à l'heure limite, et non le déséquilibre final~; exécutions dans l'enchère.
\sfield{Sources} La simulation de ce chapitre~; les règles d'enchère figurant dans l'encadré daté.
\end{strategyfile}

\begin{strategyfile}{Calage du flux de couverture}\label{strat:s1:intraday-patterns:hedging}
\sfield{Qui paie, et pourquoi} La couverture mécanique des dealers, lorsque leur exposition gamma est connue ou estimable.
\sfield{Instruments et plateformes} Contrats à terme sur indices et options.
\sfield{Signal} Une estimation du gamma net des dealers à partir des positions ouvertes sur options, combinée au mouvement de la journée.
\sfield{Dimensionnement et exécution} Ajuster la transaction de la dernière demi-heure au gamma estimé~; l'inverser lorsque le gamma est positif.
\sfield{Coûts} Comme pour le momentum de la dernière demi-heure.
\sfield{Comment elle disparaît} Estimations erronées du gamma~; changements dans l'identité des acteurs vendeurs de gamma.
\sfield{Horizon, capacité, infrastructure} Intrajournalier~; données sur le positionnement en options.
\sfield{Backtest honnête} Estimations du gamma à partir des données antérieures à chaque journée.
\sfield{Sources} Baltussen, Da, Lammers et Martens (2021) sur le lien avec la couverture en gamma.
\end{strategyfile}

\begin{strategyfile}{Retournement de l'enchère d'ouverture}\label{strat:s1:intraday-patterns:openreversal}
\sfield{Qui paie, et pourquoi} Le flux d'ordres overnight qui pousse l'ouverture, puis se retourne au cours de la journée.
\sfield{Instruments et plateformes} Les actions et l'enchère d'ouverture.
\sfield{Signal} Le \omterm{def:s1:intraday-patterns:returns}{rendement overnight}, rapporté à l'historique propre de l'action et à celui de ses pairs.
\sfield{Dimensionnement et exécution} À contre-courant du mouvement overnight depuis l'ouverture~; position nulle à la clôture.
\sfield{Coûts} Fourchettes plus larges juste après l'ouverture.
\sfield{Comment elle disparaît} Poursuite du mouvement dans la journée~; nouvelles à l'origine du mouvement.
\sfield{Horizon, capacité, infrastructure} Quelques heures~; données de l'enchère d'ouverture.
\sfield{Backtest honnête} Prix d'ouverture officiels~; aucun signal tiré de prix qui n'auraient pu être négociés.
\sfield{Sources} Lou, Polk et Skouras (2019) sur le retournement inter-périodes~; chapitre~2.
\end{strategyfile}

\section{Tutoriel~: le marché gagne son argent la nuit}\label{tut:s1:intraday-patterns:tut}

\begin{tutorial}
\textbf{Objectif.} Simuler la journée par demi-heures, avec et sans flux de couverture, et mesurer la répartition des rendements, la prévisibilité de la dernière demi-heure et le
transaction sur le déséquilibre de clôture. \textbf{État final~:} le tableau et la \cref{fig:s1:intraday-patterns:smile}.
\begin{tutsteps}
\item \textbf{La journée}~: \omterm{def:s1:intraday-patterns:returns}{rendements overnight} et par demi-heure, la poussée de couverture et le déséquilibre.
\omcode{code/firm/intraday/firm_intraday.py}{44}{59}{Modèle de journée.}\label{lst:s1:intraday-patterns:day}
\item \textbf{La régression}~: la dernière demi-heure sur le reste de la journée.
\omcode{code/firm/intraday/firm_intraday.py}{67}{74}{Prévisibilité de la dernière demi-heure.}\label{lst:s1:intraday-patterns:reg}
\item \textbf{Exécuter} \texttt{summary(0.0)}, \texttt{summary(0.05)} et \texttt{fig\_intraday.py}.
\end{tutsteps}
\textbf{Pour aller plus loin.} Faire en sorte que les dealers soient longs en gamma certains jours et vendeurs les autres, et donner à la transaction une estimation bruitée de ces jours-là~; lancer des échantillons de
cent vingt ans et tracer la distribution de la part overnight mesurée.
\end{tutorial}

\section{Construction~: la journée de trading}\label{bld:s1:intraday-patterns:intraday}

\begin{build}
\bfield{Rôle} Un modèle de la journée de trading par demi-heures, avec son \omterm{def:s1:intraday-patterns:smile}{smile du volume}, son flux de couverture et son déséquilibre de clôture, ainsi que les décompositions qui les étudient.
\bfield{Interface} \texttt{DayConfig}, \texttt{simulate\_days(n, cfg, rng)}, \texttt{decompose(open\_, close)}, \texttt{last\_half\_hour(bars, overnight)},
\texttt{volume\_curve(bars, u)}.
\bfield{Règles} Rendements espérés proportionnels à la variance, sauf imposition contraire~; la couverture et le déséquilibre agissent sur la dernière demi-heure et s'inversent
pendant la nuit.
\bfield{Tests d'acceptation} \texttt{code/firm/intraday/tests/}~: la courbe de volume et la décomposition à la main~; la répartition de la variance~; aucune prévisibilité sans
couverture, et la pente imposée avec elle.
\bfield{Pour aller plus loin} Journées par action avec enchères d'ouverture~; une estimation du gamma à partir des positions ouvertes sur options~; la bande du livre~7, chapitre~2, pour la version au niveau des ordres.
\end{build}

\begin{omsources}
\item O.~Bondarenko et D.~Muravyev, ``Market return around the clock: a puzzle'', \emph{Journal of Financial and Quantitative Analysis} 58(3), 2022 (en ligne~; numéro de 2023).
\item D.~Lou, C.~Polk et S.~Skouras, ``A tug of war: overnight versus intraday expected returns'', \emph{Journal of Financial Economics} 134(1), 2019.
\item G.~Baltussen, Z.~Da, S.~Lammers et M.~Martens, ``Hedging demand and market intraday momentum'', \emph{Journal of Financial Economics} 142(1), 2021.
\item M.~J.~Cooper, M.~T.~Cliff et H.~Gulen, ``Return differences between trading and non-trading hours: like night and day'', document de travail, 2008.
\item Nasdaq, Nasdaq Closing Cross FAQ.
\end{omsources}

\section{Exercices}

\begin{exercise}[$\star$]\label{exo:s1:intraday-patterns:1}
Avec une volatilité annuelle de 16~\%, calculer l'erreur type d'un rendement moyen sur vingt ans, puis celle de sa part overnight si un quart de la variance est overnight.
\end{exercise}

\begin{exercise}[$\star$]\label{exo:s1:intraday-patterns:2}
Si les rendements espérés sont proportionnels à la variance et qu'un quart de la variance est overnight, calculer la part d'un rendement annuel de 7~\% gagnée overnight.
\end{exercise}

\begin{exercise}[$\star$]\label{exo:s1:intraday-patterns:3}
Un déséquilibre d'un écart type pousse le cours de clôture de dix points de base, et 60~\% se retournent overnight. Calculer le gain d'une stratégie qui parie sur ce retournement, puis le gain annuel si on
l'obtenait chaque jour.
\end{exercise}

\begin{exercise}[$\star\star$]\label{exo:s1:intraday-patterns:4}
Les dealers sont vendeurs en gamma et le marché progresse de 2~\% à 15~h~30. Avec un coefficient de couverture de 0.05, calculer la poussée que le modèle ajoute à la dernière demi-heure.
\end{exercise}

\begin{exercise}[$\star\star$]\label{exo:s1:intraday-patterns:5}
Expliquer pourquoi une couverture vendeuse en gamma va dans le sens du mouvement, et une couverture longue en gamma à contre-courant.
\end{exercise}

\begin{exercise}[$\star\star$]\label{exo:s1:intraday-patterns:6}
Expliquer pourquoi vingt ans de données ne suffisent pas à établir que l'essentiel du rendement du marché est gagné overnight, alors que le résultat de Bondarenko et Muravyev est convaincant.
\end{exercise}

\begin{exercise}[$\star\star\star$]\label{exo:s1:intraday-patterns:7}
\emph{Programmation.} Exécuter \texttt{summary(0.0)} et \texttt{summary(0.05)}~; expliquer pourquoi le repli sur le déséquilibre de clôture est le même dans les deux cas, et la transaction de la dernière demi-heure
ne l'est pas.
\end{exercise}

\begin{exercise}[$\star\star\star$]\label{exo:s1:intraday-patterns:8}
\emph{Trouver la faille.} «~Les \omterm{def:s1:intraday-patterns:returns}{rendements overnight} ont été supérieurs aux rendements de séance dans notre échantillon, donc nous ne détenons le marché que la nuit et nous doublons notre
ratio de Sharpe.~»
\end{exercise}

\section{Problème~: Le marché gagne son argent la nuit}

\begin{problem}[{Problème du week-end --- la forme de la journée}]\label{pb:s1:intraday-patterns:1}
Le modèle de journée du chapitre et le constat public.

\medskip\noindent\textbf{Partie I --- La nuit et le jour.}
\begin{enumerate}
\item Définir les \omterm{def:s1:intraday-patterns:returns}{rendements overnight} et intrajournaliers.
\item Que trouvent Bondarenko et Muravyev, ainsi que Lou, Polk et Skouras~?
\item Quelle part du rendement du modèle est placée overnight, et quelle part est mesurée~?
\item Pourquoi les deux diffèrent-ils~?
\end{enumerate}

\medskip\noindent\textbf{Partie II --- La clôture.}
\begin{enumerate}[resume]
\item Définir le \omterm{def:s1:intraday-patterns:smile}{smile du volume} et donner les parts du modèle.
\item Décrire le calendrier de la confrontation de clôture.
\item Qui négocie à la clôture, et pourquoi~?
\item Que rapporte le repli sur le déséquilibre de clôture~?
\end{enumerate}

\medskip\noindent\textbf{Partie III --- Flux de couverture.}
\begin{enumerate}[resume]
\item Définir le \omterm{def:s1:intraday-patterns:hedging}{flux de couverture de fin de journée}.
\item Que trouvent Baltussen et ses coauteurs~?
\item Donner la pente de la dernière demi-heure et la transaction avec et sans couverture.
\item Que devient la configuration lorsque les dealers sont longs en gamma~?
\end{enumerate}

\medskip\noindent\textbf{Partie IV --- Le verdict.}
\begin{enumerate}[resume]
\item Énoncer le \emph{résultat nommé}~: la part du rendement cumulé gagnée overnight et la prévisibilité de la dernière demi-heure avec et sans
flux de couverture.
\item Quel comportement relève d'une prime de risque, et lequel d'une transaction de liquidité~?
\item Comment estimer le gamma des dealers~?
\item Quels coûts comptent le plus pour ces transactions~?
\item Que faut-il vérifier avant de trader la fenêtre overnight~?
\item Quelle fiche de stratégie possède la plus grande capacité~?
\item Comment ces configurations interagissent-elles avec les événements d'indices du chapitre~10~?
\item En une phrase~: qu'est-ce qui rend la clôture spéciale~?
\end{enumerate}
\end{problem}

\section{Questions d'entretien}

\begin{interviewq}[$\star$ \iqroles{trader}]\label{iq:s1:intraday-patterns:1}
Pourquoi le volume est-il le plus élevé à l'ouverture et à la clôture~?
\end{interviewq}

\begin{interviewq}[$\star\star$ \iqroles{chercheur, trader}]\label{iq:s1:intraday-patterns:2}
Qu'est-ce que le momentum intrajournalier de marché, et que pourrait-il causer~?
\end{interviewq}

\begin{interviewq}[$\star\star$ \iqroles{trader}]\label{iq:s1:intraday-patterns:3}
Comment exploiteriez-vous le déséquilibre publié de l'enchère de clôture~?
\end{interviewq}

\begin{interviewq}[$\star\star$ \iqroles{chercheur}]\label{iq:s1:intraday-patterns:4}
L'essentiel du rendement du marché semble venir overnight. Comment vérifieriez-vous que c'est réel~?
\end{interviewq}

\begin{interviewq}[$\star\star$ \iqroles{risque}]\label{iq:s1:intraday-patterns:5}
Quels risques porte une stratégie qui ne détient des positions que la nuit~?
\end{interviewq}

\begin{interviewq}[$\star\star\star$ \iqroles{chercheur}]\label{iq:s1:intraday-patterns:6}
Les dealers sont vendeurs en gamma $\Gamma$ en termes de dollars et se couvrent à la clôture. Dérivez la taille de leur transaction de couverture après un mouvement de marché $\Delta S$, et la poussée de prix
sous un impact linéaire $\lambda$ par dollar négocié.
\end{interviewq}
